% !TeX root = ../main.tex
\section{Introduction}
\label{sec:introduction}

The Video Browser Showdown (VBS) evaluates interactive video search systems
through live tasks over large collections~\cite{vbs2027cfp}. In Known-Item
Search (KIS), a participant translates a description of a previously observed
scene into search operations, inspects the returned evidence, and identifies a
valid frame or video segment~\cite{vbs2027cfp}. This process becomes more
difficult when the description contains several ordered events, because the
system must determine both what visual moments the description refers to and
where those moments occur in a candidate video.

VBS 2027 includes KIS, Ad-Hoc Video Search (AVS), and Visual Question Answering
(VQA), with most tasks using the complete Vimeo Creative Commons Collection
(V3C) together with specialized marine and surgical datasets
~\cite{vbs2027cfp}. Across its three shards, V3C contains 28,450 videos,
approximately 3,800 hours of content, and more than 4.1 million predefined
segments~\cite{vbs2027cfp,rossetto2019v3c}. The challenge also includes the
Marine Video Kit and the GynSurg collection, covering underwater footage and
laparoscopic gynecology procedures, respectively
~\cite{truong2023marine,nasirihaghighi2025gynsurg}.

Prior work on interactive CLIP search considers situations in which users
cannot formulate an effective textual description and uses intermediate
results to support query reformulation~\cite{lokoc2023queryreformulation}.
Interaction studies at VBS likewise show that query revision and result
inspection remain central to search. An analysis of two systems at VBS 2026
found iterative, high-frequency text-query reformulation followed by
result-set inspection to be the dominant strategy across most task categories
~\cite{jackl2026interaction}. A post-evaluation of VBS 2022 also showed that
the strongest initial text-to-image retrieval did not produce the highest
overall result; browsing and image-based search contributed to task
performance~\cite{schall2024interactive}. SHI builds on these observations by
making two intermediate decisions explicit for multi-event KIS: how the
request is divided into ordered events and which occurrence of each event is
selected within a candidate video.

We present SHI, an interactive video retrieval system that represents query
interpretation and selected-video alignment as two explicit hypotheses. A
\emph{Query Hypothesis} is a source-grounded ordered sequence of retrievable
events derived from the participant's original description. Before retrieval,
the participant can inspect and revise this structure through splitting,
merging, reordering, editing, or adding events. Applying an edit creates a new
revision, so retrieval remains bound to a visible query state rather than to
an implicit model-generated rewrite.

After retrieval, SHI represents a multi-event result as a complete ordered
path through one candidate video. In \emph{EventTrail}, the participant can
focus on an event and inspect temporally distinct alternative occurrences,
each shown within a complete path for all events rather than as an isolated
frame. The participant can keep the current occurrence, use an alternative, or
reject its temporal region. EventTrail converts these actions into anchor or
exclusion constraints and re-decodes the path under the updated constraints;
inspection alone does not modify the active hypothesis.

The main contributions of the submitted system are:
\begin{itemize}
    \item an inspectable, source-grounded Query Hypothesis for explicit
    correction of multi-event query structure;
    \item EventTrail, which exposes event-conditioned complete-path
    alternatives and converts user corrections into constraints for temporal
    re-decoding; and
    \item an integrated VBS workflow connecting multimodal retrieval,
    temporal alignment, hypothesis inspection, and competition-compatible
    frame submission.
\end{itemize}

The remainder of this paper describes related work, the architecture and
retrieval pipeline of SHI, the Query Hypothesis and EventTrail interactions,
and the system's evaluation and VBS workflow.
